\begin{figure}
\centering
\scalebox{0.85}{
\begin{tikzpicture}

% ==============================
% 1. FARBEN DEFINIEREN
% ==============================
\definecolor{plotblue}{RGB}{0, 0, 200}
\definecolor{plotred}{RGB}{220, 30, 30}
\definecolor{plotgray}{RGB}{140, 140, 140}
\definecolor{plotbg}{RGB}{244, 244, 252}

% ==============================
% 2. HINTERGRUND (Konvergenzbereich)
% ==============================
\fill[plotbg] (1.5, -4.8) rectangle (5.0, 4.8);

% ==============================
% 3. ACHSEN
% ==============================
\draw[thick, -{Stealth[length=3.5mm, width=2.5mm]}] (-5.2, 0) -- (5.5, 0) node[right] {\textbf{Re}$(s)$};
\draw[thick, -{Stealth[length=3.5mm, width=2.5mm]}] (0, -5.2) -- (0, 5.2) node[above] {\textbf{Im}$(s)$};

% ==============================
% 4. LINIEN & BESCHRIFTUNGEN
% ==============================
% Konvergenzabszisse (\gamma_0)
\draw[thick, plotgray, dashed] (1.5, -4.8) -- (1.5, 4.8);
\draw[thick] (1.5, 0.15) -- (1.5, -0.15) node[below right, inner sep=2pt] {$\boldsymbol{\gamma_0}$};

% Bromwich-Pfad (\gamma)
\draw[thick, plotblue, ->-=0.55] (2.5, -4.8) -- (2.5, 4.8);
\draw[thick] (2.5, 0.15) -- (2.5, -0.15) node[below right, inner sep=2pt] {$\boldsymbol{\gamma}$};

% Text im Konvergenzbereich
\node[plotblue, align=center, font=\small] at (3.75, 3) {Konvergenz-\\bereich};

% ==============================
% 5. POLE (Rote Kreuze)
% ==============================
\newcommand{\pole}[2]{
    \draw[thick, plotred, line cap=round] (#1-0.12,#2-0.12) -- (#1+0.12,#2+0.12);
    \draw[thick, plotred, line cap=round] (#1-0.12,#2+0.12) -- (#1+0.12,#2-0.12);
}

% Pole platzieren
\pole{1.5}{0}
\pole{0}{0}
\pole{0}{2}
\pole{0}{-2}
\pole{-1.2}{1.5}
\pole{-1.2}{-1.5}
\pole{-2}{0}
\pole{-3}{0}

% ==============================
% 6. LEGENDE
% ==============================
\node[draw, thick, rounded corners=3pt, fill=white, inner xsep=5pt, inner ysep=4pt, anchor=north west] at (-5.0, 4.8) {
    \renewcommand{\arraystretch}{1.2}
    \begin{tabular}{@{}c l@{}}
        
        % Symbol 1: Polstelle
        \tikz[baseline=-0.6ex]{
            \draw[thick, plotred, line cap=round] (-0.1,-0.1) -- (0.1,0.1) (-0.1,0.1) -- (0.1,-0.1);
        } & Polstelle \\
        
        % Symbol 2: Konvergenzabszisse (etwas verlängert auf -0.3 bis 0.3)
        \tikz[baseline=-0.6ex]{
            \draw[thick, plotgray, dashed] (-0.3, 0) -- (0.3, 0);
        } & Konvergenzabszisse \\
        
        % Symbol 3: Bromwich-Pfad (Mit Dekoration für perfekten Strich)
        \tikz[baseline=-0.6ex]{
            \draw[thick, plotblue, ->-=0.6] (-0.3, 0) -- (0.3, 0); % Durchgehender Strich
        } & Bromwich-Pfad \\
        
    \end{tabular}
};

\end{tikzpicture}
}
\caption{Die Bromwich-Kontur in der komplexen $s$-Ebene. Die Integrationsgerade (blau) liegt rechts von der Konvergenzabszisse $\gamma_0$ und allen Singularitäten (rote Kreuze).}
\label{laplace:fig:bromwich}
\end{figure}
